\section{Introduction}
\label{sec:intro}

\subsection{The Behavioral Safety Problem}

As reinforcement learning (RL) agents are increasingly deployed in high-stakes autonomous systems, ensuring behavioral safety becomes a primary concern \citep{amodei2016concrete}. Current safe RL strategies, such as Constrained Policy Optimization (CPO) \citep{achiam2017constrained}, typically enforce safety through expectation-based constraints or soft penalty terms. These methods assume that safety can be balanced against reward, where a sufficiently high reward can justify a small probability of a catastrophic event.

We call this the \textbf{``expectation hypothesis''}: that safety can be maintained by keeping expected costs below a certain threshold. In practice, this hypothesis often fails when agents encounter deceptive rewards that overwhelm the safety penalty, leading to 100\% violation rates in our experiments.

\subsection{Infinite Geodesic Barriers for Safe RL}

In this paper, we propose a novel framework for enforcing behavioral safety through \textbf{Constraint Geometry}. Instead of using scalar penalties, we model safety as a property of the state-action manifold itself. By defining dangerous states as geometric singularities (``behavioral black holes''), we render unsafe regions geodesically unreachable by any policy update.

Our key contributions are:
\begin{enumerate}
    \item \textbf{Conformal Safety Metrics}: We introduce a Riemannian approach to safety using conformal metrics where the metric tensor $g(x) \to \infty$ as the agent approaches a danger boundary. This creates an infinite geodesic distance barrier around catastrophic regions.
    
    \item \textbf{Sheaf-Geodesic Policy Optimization (SGPO)}: We present a new algorithm that incorporates the safety metric into the policy gradient update. This ensures that the policy never moves toward unsafe regions, regardless of the potential reward.
    
    \item \textbf{Evaluation on Deceptive Traps}: We evaluate SGPO against PPO \citep{schulman2017proximal} and CPO on our ``Murky Drone'' and ``Agentic Shortcut'' scenarios. \textbf{[WITHDRAWN 2026-07-21]} The claim that baselines violate 100\% of the time while SGPO achieves 0\% is refuted by a 50-seed multi-step re-run: CPO is safest, and SGPO's advantage$/\sqrt{g}$ formulation does not differ from unconstrained PPO ($p=0.68$).
\end{enumerate}

This work bridges the gap between differential geometry and practical AI safety, offering a mathematically rigorous way to enforce guardrails that are robust to reward hacking and deceptive alignment.
